\documentclass{article}

\usepackage{amsmath,amssymb}
\usepackage{xurl}
\usepackage[hidelinks]{hyperref}
\setlength{\emergencystretch}{2em}

% Shared commands for notes in docs/paper-gaps/.

\DeclareUnicodeCharacter{2080}{\ifmmode{}_{0}\else\textsubscript{0}\fi}
\DeclareUnicodeCharacter{2081}{\ifmmode{}_{1}\else\textsubscript{1}\fi}
\DeclareUnicodeCharacter{2082}{\ifmmode{}_{2}\else\textsubscript{2}\fi}
\DeclareUnicodeCharacter{2099}{\ifmmode{}_{n}\else\textsubscript{n}\fi}

\newcommand{\Lean}{\textsc{Lean}}
\ExplSyntaxOn
\NewDocumentCommand{\leanid}{m}{
  \ifmmode
    \textsf{\detokenize{#1}}
  \else
    \group_begin:
    \urlstyle{sf}
    \tl_set:Nn \l_tmpa_tl {#1}
    \tl_replace_all:Nnn \l_tmpa_tl {\_} {_}
    \exp_args:NV \path \l_tmpa_tl
    \group_end:
  \fi
}
\ExplSyntaxOff
\providecommand{\tr}{\operatorname{tr}}
\newcommand{\tnleanIssueBase}{https://github.com/LionSR/TNLean/issues/}
\newcommand{\tnleanPrBase}{https://github.com/LionSR/TNLean/pull/}
\newcommand{\tnleanDocBase}{https://sirui-lu.com/TNLean/docs/}
\newcommand{\ghissue}[1]{\href{\tnleanIssueBase#1}{GitHub issue \##1}}
\newcommand{\ghpr}[1]{\href{\tnleanPrBase#1}{PR \##1}}
\newcommand{\doclink}[1]{\href{\tnleanDocBase#1.html}{\path{#1}}}

% Dirac notation and the identity operator, as in blueprint/src/macros/common.tex.
% \providecommand keeps note-local definitions authoritative.
\usepackage{dsfont}
\providecommand{\ket}[1]{|#1\rangle}
\providecommand{\bra}[1]{\langle #1|}
\providecommand{\braket}[2]{\langle #1 | #2 \rangle}
\providecommand{\ketbra}[2]{|#1\rangle\!\langle #2|}
\providecommand{\Id}{\mathds{1}}
\providecommand{\one}{\mathds{1}}
\providecommand{\C}{\mathbb{C}}
\providecommand{\MN}[1]{M_{#1}(\C)}
\providecommand{\MD}{M_D(\C)}

% Verdict marker: \gapnote{<kind>}{<status>}. Kind is the policy
% classification (clarification, local-correction, scope-restriction,
% unfaithful, false-source, open-gap); status is open, wip (elimination
% underway), resolved, or historical. Machine-read by the site index;
% prints a verdict line.
\newcommand{\gapnote}[2]{\par\noindent\textbf{Verdict:} #1 (#2).\par}


% Fallbacks keep this note compilable when a different command.tex is found first.
\providecommand{\leanid}[1]{\textsf{\detokenize{#1}}}
\providecommand{\gapnote}[2]{\par\noindent\textbf{Verdict:} #1 (#2).\par}

\title{The Logarithmic Factor in the Matrix Product State Approximation Quoted by Landau, Vazirani and Vidick}
\author{}
\date{2026-11-03}

\begin{document}
\maketitle
\gapnote{local-correction}{resolved}

\section{The source statement}

Landau, Vazirani and Vidick~\cite{Landau2015Polynomial},
arXiv:1307.5143v1, \path{algo53.tex}, lines~325--328, state the following
lemma for a nearest-neighbour chain of $n$ qudits of dimension $d$ with
$0\le H_i\le1$, unique ground state $\Gamma$ and gap $\varepsilon$.
\begin{quote}
    For any constant $c>0$ there is a constant $C\ge1$ such that for every
    $n$ there is a vector $v$ with Schmidt rank bounded by
    $\exp(C(\ln n)^{3/4}\varepsilon^{-1/4})$ across every cut such that
    $\lvert\langle\Gamma,v\rangle\rvert\ge1-n^{-c}$.
\end{quote}
No proof is given. The lemma is quoted from Arad, Kitaev, Landau and
Vazirani~\cite{Arad2013AreaLaw}, and the source adds (line~325) that only a
polynomial bound on the Schmidt rank is needed for its algorithm.

\section{The cited bound}

The cited work, arXiv:1301.1162v1, \path{1dfrustrated34.tex}, treats $d$ as a
constant and states the bond dimension as
$B=e^{\tilde O(\log^{3/4}n/\varepsilon^{1/4})}$, both in its introduction
(line~176) and in the derivation (lines~1503--1517). There the rank of the
approximate ground-state projector is $\tilde O(\exp(\ell/s+s))$ with
$\ell=O(\log n\sqrt{(s+\log n)/\varepsilon})$, and the choice
$s=\log^{3/4}n/\varepsilon^{1/4}$ balances the two terms. The tilde hides
factors logarithmic in the exponent; the quotation in the source omits it.

\section{What is formalized}

For a nearest-neighbour chain of $L$ sites of dimension $q$ with terms
$0\le h_j\le1$, a unit ground vector $\Omega$ of energy $\mu$ and an actual
first excited energy $\mu+g$ with $g>0$, and for $0<\epsilon\le1$, put
\begin{align*}
    u=1+\log q+\log\frac1g,\qquad
    z=\log\frac{2048(L+1)}{\epsilon^2}.
\end{align*}
There are universal constants $C>0$ and a bond dimension $D$ with
\begin{align}
    \log D\le C\Bigl[\frac{u^3}{g}
        +\Bigl(\frac{z^3}{g}\Bigr)^{1/4}\bigl(u+\log(z+1)\bigr)\Bigr]
    \label{eq:lvv13_mpsapprox_formal}
\end{align}
and a normalized open-boundary matrix product state $\Phi$ of bond dimension
at most $D$ with $\lVert\Phi-\Omega\rVert\le\epsilon$.\footnote{%
    \leanid{TNLean.MPS.AreaLaw.exists_aklv_sharp_physical_sequential_mps_approximation}
    in \doclink{TNLean/MPS/AreaLaw/AKLV/SharpMPSApproximation}.}
The overlap then satisfies
$\lvert\langle\Omega,\Phi\rangle\rvert\ge1-\epsilon^2/2$.

The theorem requires $\mu+g$ to lie in the spectrum, so that $g$ is the gap
itself and not only a lower bound for it. This is the source's setting, where
$\varepsilon=\varepsilon_1-\varepsilon_0$ is the difference of the two lowest
eigenvalues (line~283). The local support and normalization conventions of the
underlying construction are recorded in
\path{docs/paper-gaps/aklv13_two_block_robustness.tex}.

\section{Comparison}

Take $\epsilon=n^{-c/2}$ with $n=L$. Then $z=O(\log n)$, and for fixed $q$ and
$g$ the bound (\ref{eq:lvv13_mpsapprox_formal}) reads
\begin{align*}
    \log D=O\bigl(g^{-1/4}(\log n)^{3/4}\log\log n\bigr),
\end{align*}
with overlap at least $1-n^{-c}/2$. This is the cited bound
$e^{\tilde O(\log^{3/4}n/\varepsilon^{1/4})}$, with the hidden factor
written out as $\log\log n$. It is smaller than any positive power of $n$,
so it is sub-polynomial. The printed lemma, without the factor
$\log\log n$, is not formalized and is not claimed by the cited work. The
gap enters (\ref{eq:lvv13_mpsapprox_formal}) also through $u$ and the term
$u^3/g$, which are independent of $n$.

The algorithm of the source uses only a polynomial bound on the Schmidt
rank, so no later step depends on the factor. The viable-set construction
uses the polynomial bound of the library's one-dimensional area law.

\section{Verdict}

A local correction, resolved: Lemma \path{l:mpsapprox} is read with the
$\tilde O$ of its cited source, and in that form it is formalized under the
source's hypotheses.

\bibliographystyle{plain}
\bibliography{references}

\end{document}
